\documentclass{article} %210 mm × 297 mm
\usepackage[utf8]{inputenc}
\usepackage[a4paper,left=1.5cm, right=1.5cm, top=2cm, bottom=2cm]{geometry}
\usepackage{graphicx}
%\usepackage{blindtext}
\usepackage{multirow}
\usepackage{mhchem}
\usepackage{url}
\usepackage{subfigure}
\usepackage{amsfonts,latexsym} % para tener disponibilidad de diversos simbolos
\usepackage{enumerate}
%\usepackage{booktabs}
\usepackage{float}
%\usepackage{threeparttable}
\usepackage{array,colortbl}
%\usepackage{ifpdf}
%\usepackage{rotating}
\usepackage{cite}
\usepackage{stfloats}
\usepackage{url}
\usepackage{listings}
\usepackage{fancyhdr}

\usepackage{color}
%\usepackage{hyperref}
%\usepackage{wrapfig}
\usepackage{array}
%\usepackage{multirow}
%\usepackage{adjustbox}
%\usepackage{nccmath}
\usepackage[dvipsnames]{xcolor}


\newcommand{\titulo}{Abgabe 1; MfN II.}
\newcommand{\nombre}{Liliana Sanfilippo}
\newcommand{\FCE}{\textbf{SoSe 2025}}

 

\usepackage{fancyhdr}
\pagestyle{fancy}
\fancyhead{}
\fancyfoot{}
\fancyhead[LO]{\small\nombre}
\fancyhead[RE]{\small\titulo}
\fancyhead[LO]{\small\FCE}
\fancyfoot[RO,LE] {\thepage}

\setcounter{page}{1} %Número de la página, la editorial lo cambiará



\usepackage[onehalfspacing]{setspace}
\begin{document}


\begin{tabular}{p{12cm}} 
{{\Large\bf\titulo}}\\
\vspace{0.2cm}\\
 {\large\nombre}\\
\end{tabular}
\vspace{0.2cm}\\
\section{Aufgabe 1.1}
(Regel von Sarrus)
\begin{align*}
det 
\begin{pmatrix}
1 & 2 & 3 \\
2 & 5 & 1 \\
2 & 7 & 9
\end{pmatrix} & = 
(1 \cdot 5 \cdot 9) + (2 \cdot 1 \cdot 2) + (3 \cdot 2 \cdot 7) - (3\cdot5\cdot2) - (1 \cdot 1 \cdot 7) - (2 \cdot 2 \cdot 9) \\
&  = 45 + 4 + 42 - 30 - 7 - 36 \\
& = 18
\end{align*} 

\section{Aufgabe 1.2}
(Regel von Sarrus)
\begin{align*}
det 
\begin{pmatrix}
x & 1 & 1 \\
1 & x & 1 \\
1 & 1 & x
\end{pmatrix} & = x^3 + 1 + 1 - x - x - x \\
& = x^3 +2 -3x  \\
& = x^3 - 3x +2 \\
& = x^3 +2x^2 -2x^2 - 4x +x +2 \\
& = (x^2 + 2x +1) \cdot (x+2) \\
& = (x-1)^2 \cdot (x+2)
\end{align*} 

\section{Aufgabe 1.3}
(Laplacescher Entwicklungssatz) 
\[
A = \begin{pmatrix}
\sin \alpha & \cos \alpha & a \cdot \sin \alpha & b \cdot \cos \alpha & ab \\
- \cos \alpha & \sin \alpha & -a^2 \sin \alpha & b^2 \cos \alpha & a^2 b^2 \\
0 & 0 & 1 & a^2 & b^2 \\
0 & 0 & 0 & a & b \\
0 & 0 & 0 & -b & -a
\end{pmatrix}
\]
Entwicklung nach erster Spalte ($j = 1$). 
\begin{align*}
    det(A) & = \sum^5_{i =1} a_{i1} \cdot (-1)^{1+i} \cdot det(A'_{i1})  \\
    & = a_{11} \cdot det(A'_{11}) + \sum^5_{i =2} a_{i1} \cdot (-1)^{1+i} \cdot det(A'_{i1}) \\
    & = a_{11} \cdot det(A'_{11}) + a_{21} \cdot (-1)\cdot det(A'_{21})  \\
    & = a_{11} \cdot det(A'_{11}) - a_{21} \cdot det(A'_{21})  
\end{align*}
Da $a_{i3} - a_{i5}$ Null sind.
\[
A'_{11} = \begin{pmatrix}
 \sin \alpha & -a^2 \sin \alpha & b^2 \cos \alpha & a^2 b^2 \\
 0 & 1 & a^2 & b^2 \\
 0 & 0 & a & b \\
0 & 0 & -b & -a
\end{pmatrix}
\]
\[
A'_{21} = \begin{pmatrix}
 \cos \alpha & a \cdot \sin \alpha & b \cdot \cos \alpha & ab \\
 0 & 1 & a^2 & b^2 \\
 0 & 0 & a & b \\
 0 & 0 & -b & -a
\end{pmatrix}
\]

\subsection*{det($A'_{11}$)}
Laplace-Entwiklung nach der ersten Spalte. Da $a_{i2} - a_{i1}$ alle Null gilt: 
\begin{align*}
 det(A'_{11}) & = \sum^3_{i =1} a_{i1} \cdot (-1)^{1+i} \cdot det(A''_{i1})  = \sin \alpha \cdot det(A''_{11})  \\ \\
A''_{11} &  = \begin{pmatrix}
1 & a^2 & b^2 \\
 0 & a & b \\
 0 & -b & -a
\end{pmatrix} \\ \\
    det(A''_{11})  & =  \sum^2_{i =1} a_{i1} \cdot (-1)^{1+i} \cdot det(A'''_{i1})  \\
    & =  1 \cdot det(A'''_{11})  \\ \\
    A'''_{11}  & = \begin{pmatrix}
  a & b \\
  -b & -a
\end{pmatrix} \\ \\
det( A'''_{11}) & = -a^2 + b^2 \\ \\
Also: & \\
 det(A'_{11}) & = \sin \alpha \cdot (1 \cdot (-a^2 + b^2))  \\
 & = \sin \alpha \cdot(-a^2 + b^2) \\
 & = -a^2 \cdot \sin \alpha  + b^2 \cdot \sin \alpha
\end{align*}

\subsection*{det($A'_{21}$)}
Laplace-Entwiklung nach der ersten Spalte. Da $a_{i2} - a_{i1}$ alle Null gilt: 
\begin{align*}
     det(A'_{21}) & = \sum^3_{i =1} a_{i1} \cdot (-1)^{1+i} \cdot det(A''_{i1})  = a_{11} \cdot det(A''_{11})  \\ \\
     A''_{11} &  = \begin{pmatrix}
1 & a^2 & b^2 \\
 0 & a & b \\
 0 & -b & -a
\end{pmatrix} \\ \\
det(A''_{11})  & = -a^2 + b^2 \\ \\
Also \\
 det(A'_{21}) & = \cos \alpha \cdot(-a^2 + b^2) \\
 & = -a^2 \cdot \cos \alpha  + b^2 \cdot \cos \alpha 
\end{align*}


Also:
\begin{align*}
    det(A) & = a_{11} \cdot det(A'_{11}) - a_{21} \cdot det(A'_{21})   \\
    & = \sin \alpha \cdot (-a^2 \cdot \sin \alpha  + b^2 \cdot \sin \alpha) - \cos \alpha (-a^2 \cdot \cos \alpha  + b^2 \cdot \cos \alpha)  \\
    & = -a^2 \sin^2 \alpha  + b^2 \sin^2 \alpha + a^2 \cos^2 \alpha -b^2 \cos^2 \alpha
\end{align*}

Frage: Hätte man das an irgendeiner Stelle schöner auflösen können? 

\section{Aufgabe 1.4}
\subsection*{Induktionsanfang}
Sei \(n = 2\) und zeige $A_{n+1} = \frac{(n+2+1)!}{4} = \frac{(n+3)!}{4}$
\begin{align*}
A_2 & = \begin{pmatrix}
1 & 2 \\
-2 & 2
\end{pmatrix}   \\ \\
det (A_2 ) & = (1 \cdot 2) - (2 \cdot (-2)) \\
& = 2 + 4 \\
& = 6 \\
& = \frac{24}{4} \\
& = \frac{1\cdot 2 \cdot 3 \cdot 4 }{4} \\
& = \frac{4!}{4} \\
& = \frac{(2+2)!}{4}
\end{align*}
Die Aussage gilt also für ein $A_n$. 
\subsection*{Induktionsschritt}
\(n \Rightarrow n+1\)
\begin{align*}
    A_{n+1} & =
\begin{pmatrix}
    1 & 2 & 2 & \cdots & 2 & 2 \\
   -2 & 2 & \ddots & 2 & 2 & 2\\
   -2 & -2 & 3 & 2 & 2 & 2 \\
   -2 & -2 & \ddots & n-1 & 2 & 2\\
   -2 & -2 & \ddots & -2 & n & 2 \\
   -2 & -2 & -2 & \cdots & -2 & n+1
\end{pmatrix} \\ \\
det(A_{n+1}) & = \sum^n_{i=1} a_{i1} \cdot (-1)^{1+i} \cdot det(A'_{i1}) \\
& = 1 \cdot (-1)^{2} \cdot det(A'_{11}) + \sum^n_{i=2} a_{i1} \cdot (-1)^{1+i} \cdot det(A'_{i1}) \\
& = 1 \cdot (-1)^{2} \cdot det(A'_{11}) + \sum^n_{i=2} (-2) \cdot (-1)^{1+i} \cdot det(A'_{i1}) \\
& =  det(A'_{11}) + (-2) \cdot \sum^n_{i=2} (-1)^{1+i} \cdot det(A'_{i1}) \\
%& =  \frac{24}{4} + (-2) \cdot \sum^n_{i=2} (-1)^{1+i} \cdot det(A'_{i1}) \\
\end{align*}
Dabei ist $ \sum^n_{i=2} (-2) \cdot (-1)^{1+i} \cdot det(A'_{i1})$ die Laplace-Entwicklung der Matrix nach der ersten Spalte. Man kann die Matrix folgendermaßen umstellen, ohen Zeilen zu vertauschen, sodass die Determinante gleich bleibt: 

%\subsection{Nicht-Induktion zur Überprüfung}
\[
\begin{pmatrix}
    1 & 2 & \cdots & 2 & 2  \\
   -2 & 2 & \cdots & 2 & 2 \\
   -2 & -2 & 3 & \ddots & 2 \\
   -2 & -2 & \ddots & n & 2 \\
   -2 & -2 & \cdots & -2 & n+1 
\end{pmatrix}
\]
 \text{Zeilen 2 bis n mit der dem zweifachen der ersten Zeile addieren}
\[
\begin{pmatrix}
    1 & 2 & \cdots & 2 & 2  \\
   0 & 6 & \cdots & 6 & 6 \\
   0 & 2 & 7 & \ddots & 6 \\
   0 & 2 & \ddots & n+4 & 6 \\
   0 & 2 & \cdots & 2 & n+5 
\end{pmatrix}
\]
Die 3 als Faktor aus der zweiten Zeile herausziehen.
\[
3 \cdot \begin{pmatrix}
    1 & 2 & \cdots & 2 & 2  \\
   0 & 2 & \cdots & 2 & 2 \\
   0 & 2 & 7 & \ddots & 6 \\
   0 & 2 & \ddots & n+4 & 6 \\
   0 & 2 & \cdots & 2 & n+5 
\end{pmatrix}
\]
 \text{Zeilen 3 bis n mit der dem -1-fachen der 2-ten Zeile addieren}
 \[
3 \cdot \begin{pmatrix}
    1 & 2 & \cdots & 2 & 2  \\
   0 & 2 & \cdots & 2 & 2 \\
   0 & 0 & 5 & \ddots & 4 \\
   0 & 0 & \ddots & n+2 & 4 \\
   0 & 0 & \cdots & 0 & n+3 
\end{pmatrix}
\]
Determinante wurde nicht verändert, da keine Zeilenvertauschungen stattgefunden haben...

\begin{align*}
    det\ 3 \cdot \begin{pmatrix}
    1 & 2 & \cdots & 2 & 2  \\
   0 & 2 & \cdots & 2 & 2 \\
   0 & 0 & 5 & \ddots & 4 \\
   0 & 0 & \ddots & n+2 & 4 \\
   0 & 0 & \cdots & 0 & n+3 
\end{pmatrix} & = 1 \cdot 2 \cdot 5 \cdot 6 \cdots (n+2) \cdot (n+3) \\
   det  \begin{pmatrix}
    1 & 2 & \cdots & 2 & 2  \\
   0 & 2 & \cdots & 2 & 2 \\
   0 & 0 & 5 & \ddots & 4 \\
   0 & 0 & \ddots & n+2 & 4 \\
   0 & 0 & \cdots & 0 & n+3 
\end{pmatrix}   & = 1 \cdot 2  \cdot 3 \cdot 5 \cdot 6 \cdots (n+2) \cdot (n+3) \\
& = \frac{(n+3)!}{4}
\end{align*}
So funktioniert die Determinante für jedes $A'_{i1}, \quad i>1$, weshalb sie sich aufgrund der wechselnden Vorzeichen gegenseitig aufheben und nur $det(A'_{11}) = \frac{24}{4}$ übrig bleibt. 




\end{document}